% SOURCE: Chinese PDF viewer pp.156--159, exactly repeated on pp.163--166.
\subsection*{4) Commutation relations of creation and annihilation operators}
\textbf{A. First consider the relations between $\op a_i^\dagger$ and $\op a_j^\dagger$ (or between $\op a_i$ and $\op a_j$).}

For any state $\ket\psi$, the states $\op a_i^\dagger\op a_j^\dagger\ket\psi$ and $\op a_j^\dagger\op a_i^\dagger\ket\psi$ are the same state, and can differ only by a constant:
\[\op a_i^\dagger\op a_j^\dagger\ket\psi
=\lambda\op a_j^\dagger\op a_i^\dagger\ket\psi.\]
Consequently,
\begin{align*}
0&=(\op a_i^\dagger\op a_j^\dagger-\lambda\op a_j^\dagger\op a_i^\dagger)\ket\psi\\
&=\sum_{k,l}S_{ki}S_{lj}
(\op b_k^\dagger\op b_l^\dagger-\lambda\op b_l^\dagger\op b_k^\dagger)\ket\psi\\
&=\sum_{l,k}S_{li}S_{kj}
(\op b_l^\dagger\op b_k^\dagger-\lambda\op b_k^\dagger\op b_l^\dagger)\ket\psi.
\end{align*}
By the independence of $\ket\psi$ and the representation-transformation matrix elements $S_{ji}=\braket{b_j}{a_i}$ (that is, the independence of the basis vectors),
\begin{align*}
\op a_i^\dagger\op a_j^\dagger-\lambda\op a_j^\dagger\op a_i^\dagger&=0,\\
\op b_k^\dagger\op b_l^\dagger-\lambda\op b_l^\dagger\op b_k^\dagger&=0,\\
\op b_l^\dagger\op b_k^\dagger-\lambda\op b_k^\dagger\op b_l^\dagger&=0.
\end{align*}
The last two relations lead to
\[(1-\lambda^2)\op b_k^\dagger\op b_l^\dagger=0,\qquad\lambda=\pm1.\]
For $\lambda=+1$ the first relation gives
\[\op a_i^\dagger\op a_j^\dagger-\op a_j^\dagger\op a_i^\dagger=0,\qquad
\op a_i\op a_j-\op a_j\op a_i=0.\]
For $\lambda=-1$, it gives
\[\op a_i^\dagger\op a_j^\dagger+\op a_j^\dagger\op a_i^\dagger=0,\qquad
\op a_i\op a_j+\op a_j\op a_i=0.\]

\textbf{B. Next consider the relations between $\op a_i$ and $\op a_j^\dagger$.}

First take $i\ne j$. The vectors $\op a_i\op a_j^\dagger\ket\psi$ and $\op a_j^\dagger\op a_i\ket\psi$ describe the same state. Therefore,
\[0=(\op a_i\op a_j^\dagger-\mu\op a_j^\dagger\op a_i)\ket\psi,\]
so that
\[\op a_i\op a_j^\dagger-\mu\op a_j^\dagger\op a_i=0.\]
Under a transformation of representation,
\[0=\sum_{k,l}S_{ki}^*S_{lj}
(\op b_k\op b_l^\dagger-\mu\op b_l^\dagger\op b_k)\ket\psi.\]
If $k\ne l$, independence of the transformation-matrix elements implies
\[\op b_k\op b_l^\dagger-\mu\op b_l^\dagger\op b_k=0.\]
If $k=l$, we have
\[0=\sum_k S_{ki}^*S_{kj}
(\op b_k\op b_k^\dagger-\mu\op b_k^\dagger\op b_k)\ket\psi.\]
Only if
\[\op b_k\op b_k^\dagger-\mu\op b_k^\dagger\op b_k=\op A,\]
where $\op A$ is independent of $k$, can we ensure that
\begin{align*}
\sum_kS_{ki}^*S_{kj}
(\op b_k\op b_k^\dagger-\mu\op b_k^\dagger\op b_k)\ket\psi
&=\sum_kS_{ki}^*S_{kj}\op A\ket\psi\\
&=\delta_{ij}\op A\ket\psi=0\qquad(i\ne j).
\end{align*}
Now take $i=j$. Then
\begin{align*}
\op a_i\op a_j^\dagger-\mu\op a_j^\dagger\op a_i
&=\sum_{k,l}S_{ki}^*S_{lj}
(\op b_k\op b_l^\dagger-\mu\op b_l^\dagger\op b_k)\\
&=\sum_kS_{ki}^*S_{kj}(\op b_k\op b_k^\dagger-\mu\op b_k^\dagger\op b_k)\\
&=\sum_kS_{ki}^*S_{kj}\op A=\op A\qquad(i=j).
\end{align*}
Combining $i\ne j$ and $i=j$,
\[\op a_i\op a_j^\dagger-\mu\op a_j^\dagger\op a_i=\op A\delta_{ij}.\]
\textbf{Note:} The relation between $\op a_i$ and $\op a_j^\dagger$ differs from the relation between two creation operators (or between two annihilation operators) because of the unitarity of the transformation matrix,
\[\sum_kS_{ki}^*S_{kj}=\delta_{ij}.\]

\textbf{C. Determine $\mu$ and $\op A$.}

The total number of particles is independent of representation:
\[N=\sum_i n_i=\sum_i\widetilde n_i,\qquad
\op N=\sum_i\op n_i=\sum_i\op{\widetilde n}_i.\]
Since both $\sum_i\op a_i^\dagger\op a_i$ and $\sum_i\op a_i\op a_i^\dagger$ are independent of representation, write
\[\op N=x\sum_i\op a_i^\dagger\op a_i
+y\sum_i\op a_i\op a_i^\dagger+Z,\]
\[\op n_i=x\op a_i^\dagger\op a_i+y\op a_i\op a_i^\dagger+z,\]
where $x,y,z$ are constants. As
\[\op n_i\ket0=0,\qquad
(x\op a_i^\dagger\op a_i+y\op a_i\op a_i^\dagger+z)\ket0=0,\]
and
\[\op a_i\ket0=0,\qquad
\op a_i\op a_i^\dagger\ket0=\op a_i\ket{a_i}=\ket0,\]
we have
\[(y+z)\ket0=0,\qquad z=-y,\]
so that
\[\op n_i=x\op a_i^\dagger\op a_i+y(\op a_i\op a_i^\dagger-1).\]
Again consider $i\ne j$. In order to determine $\mu$, use
\[\op a_i\op a_j^\dagger-\mu\op a_j^\dagger\op a_i=0.\]
For the occupation-number state $\ket{\ldots,n_i,\ldots,n_j,\ldots}$,
\begin{align*}
\op n_i\op a_j\ket{\ldots,n_i,\ldots,n_j,\ldots}
&=c_j n_i\ket{\ldots,n_i,\ldots,n_j-1,\ldots},\\
\op a_j\op n_i\ket{\ldots,n_i,\ldots,n_j,\ldots}
&=n_ic_j\ket{\ldots,n_i,\ldots,n_j-1,\ldots}.
\end{align*}
Hence
\[[\op n_i,\op a_j]=0,\]
which yields
\[x(\op a_i^\dagger\op a_i\op a_j-\op a_j\op a_i^\dagger\op a_i)
+y(\op a_i\op a_i^\dagger\op a_j-\op a_j\op a_i\op a_i^\dagger)=0.\]
For $\op a_i\op a_j-\op a_j\op a_i=0$, this becomes
\[x(\op a_i^\dagger\op a_j-\op a_j\op a_i^\dagger)\op a_i
+y\op a_i(\op a_i^\dagger\op a_j-\op a_j\op a_i^\dagger)=0,\]
so
\[\op a_j\op a_i^\dagger-\op a_i^\dagger\op a_j=0.\]
Likewise, if $\op a_i\op a_j+\op a_j\op a_i=0$, we obtain
\[\op a_j\op a_i^\dagger+\op a_i^\dagger\op a_j=0.\]
Thus for $i\ne j$, the sign $\mu=\pm1$ is the same as the corresponding sign $\lambda=\pm1$.

For $i=j$, the relation is
\[\op a_i\op a_i^\dagger-\mu\op a_i^\dagger\op a_i=\op A.\]
From its action on the vacuum,
\[(\op a_i\op a_i^\dagger-\mu\op a_i^\dagger\op a_i)\ket0=\op A\ket0,\]
we obtain
\[\ket0=\op A\ket0,\qquad\op A=1.\]
Combining $i\ne j$ and $i=j$ gives
\[\boxed{\op a_i\op a_j^\dagger-\mu\op a_j^\dagger\op a_i=\delta_{ij}.}\]

\textbf{D. Summary of commutation relations.}
There are precisely two sets of relations:
\begin{align*}
\text{Bosons:}\qquad
[\op a_i,\op a_j]&=0,&
[\op a_i^\dagger,\op a_j^\dagger]&=0,&
[\op a_i,\op a_j^\dagger]&=\delta_{ij};\\
\text{Fermions:}\qquad
\{\op a_i,\op a_j\}&=0,&
\{\op a_i^\dagger,\op a_j^\dagger\}&=0,&
\{\op a_i,\op a_j^\dagger\}&=\delta_{ij}.
\end{align*}
They correspond to the two kinds of identical particles: bosons and fermions.

\subsection*{5) Pauli exclusion principle}
For a system of fermions,
\[\op a_i^\dagger\op a_i^\dagger=0,\qquad
\op a_i^\dagger\op a_i^\dagger\ket0=0,\]
\[\op a_i^\dagger\ket{0,\ldots,n_i=1,\ldots}=0,\qquad
\ket{0,\ldots,n_i=2,\ldots}=0.\]
Thus two or more fermions cannot occupy the same single-particle state. For fermions, $n_i=0,1$.

\subsection*{6) Particle-number operator}
We have
\[\op n_i=x\op a_i^\dagger\op a_i+y(\op a_i\op a_i^\dagger-1).\]
Because
\[\op a_i\op a_i^\dagger=1\pm\op a_i^\dagger\op a_i,\]
we obtain
\[\op n_i=(x\pm y)\op a_i^\dagger\op a_i.\]
Since
\[\op n_i\ket{a_i}=\ket{a_i},\]
\[(x\pm y)\op a_i^\dagger\op a_i\ket{a_i}
=(x\pm y)\op a_i^\dagger\ket0=(x\pm y)\ket{a_i}=\ket{a_i},\]
we find
\[x\pm y=1,\qquad\boxed{\op n_i=\op a_i^\dagger\op a_i}.\]